\chapter{Interferencja i doświadczenie z dwiema szczelinami}

\section{Superpozycja amplitud, nie prawdopodobieństw}
\section{Dwie szczeliny jako model interferencji}
\section{Pojedyncza szczelina jako granica}
\section{Trzy szczeliny i więcej szczelin}
\section{Szczeliny punktowe w 2D}
\section{Prążki interferencyjne jako mapa \texorpdfstring{\(|\psi(x,y)|^2\)}{|psi(x,y)|^2}}
\section{Co niszczy interferencję?}
\section{Interpretacja: fala prawdopodobieństwa a pojedyncze zdarzenia}
